\ifdefined\ismainfile
\else
    \documentclass[11pt]{article}
    
    \makeatletter
\@ifclassloaded{revtex4-2}{}{\usepackage[margin=1in]{geometry}}
\makeatother
\usepackage{amsmath,amsthm,amsfonts,amssymb,mathrsfs}
\usepackage{graphicx}
\usepackage{cancel}
\usepackage{bm}
\usepackage{tikz}
\usepackage{csquotes}
\usepackage{tikz-3dplot}
\usepackage{physics}
\usepackage{dsfont}
\usepackage{float}
\usepackage{xcolor}
\usepackage{mathtools}
\usepackage{tcolorbox}
\makeatletter
\@ifclassloaded{revtex4-2}{\let\affil\affiliation
}{\usepackage{authblk}}
\makeatother

\tcbuselibrary{breakable}
\usepackage{keytheorems}
\newkeytheorem{definition}[parent=section, tcolorbox-no-titlebar={breakable, colback=blue!10}]
\newkeytheorem{example,lemma,conjecture,claim,theorem,proposition,corollary,fact,assumption,construction}[sibling=definition, tcolorbox-no-titlebar={breakable, colback=blue!10}]
\theoremstyle{definition}
\newtheorem{protocol}{Protocol}[section]

\newkeytheorem{remark}[parent=section, tcolorbox-no-titlebar={breakable}]
\allowdisplaybreaks[4]
\interdisplaylinepenalty=0
\makeatletter
\renewcommand{\log}{\mathop{\operator@font log}\nolimits}
\makeatother
\IfFileExists{bbm.sty}{\usepackage{bbm}}{\providecommand{\mathbbm}{\mathds}}
\usepackage{stmaryrd}
\usepackage{braket}
\makeatletter
\@ifpackageloaded{natbib}{}{\usepackage[numbers,sort&compress]{natbib}}
\makeatother
\usepackage{hyperref}
\usepackage[ruled,noline,longend]{algorithm2e}
\SetAlgoNoLine
\AtBeginDocument{\SetAlgoNoLine}
\SetAlgoLongEnd
\RestyleAlgo{boxruled}
\LinesNumbered
\definecolor{algcommentgray}{gray}{0.35}
\newcommand{\AlgCommentFont}[1]{\textcolor{algcommentgray}{#1}}
\SetCommentSty{AlgCommentFont}
\SetKwComment{tcp}{{\color{algcommentgray}$\blacktriangleright$}\ }{}
\usepackage{tcolorbox}
\tcbuselibrary{breakable}
\renewcommand{\topfraction}{0.9}
\renewcommand{\bottomfraction}{0.8}
\renewcommand{\textfraction}{0.07}
\renewcommand{\floatpagefraction}{0.85}
\setcounter{topnumber}{3}
\setcounter{totalnumber}{5}
\newenvironment{longalgorithm}[2][]{\begin{tcolorbox}[breakable, sharp corners, boxsep=0pt,
      boxrule=\algoheightrule, colframe=black, colback=white,
      left=18pt, right=4pt, top=5pt, bottom=5pt, #1]\refstepcounter{algocf}\noindent\hspace*{-14pt}{\textbf{\algorithmcfname~\thealgocf:} #2}\par
  \addcontentsline{loa}{algocf}{\protect\numberline{\thealgocf}#2}\par\nobreak\vspace{2pt}\nointerlineskip
  \moveleft18pt\vbox{\hrule width \dimexpr\linewidth+22pt\relax height \algoheightrule}\vspace{3pt}\setcounter{AlgoLine}{0}\RestyleAlgo{plain}}{\end{tcolorbox}}
\DontPrintSemicolon
\SetKwInput{KwIn}{Require}
\SetKwInput{KwOut}{Ensure}
\SetKwFor{For}{for}{do}{end for}
\SetKwFor{ForEach}{for each}{do}{end for}
\SetKwFor{While}{while}{do}{end while}
\makeatletter
\newenvironment{breakablealgorithm}{\par\addvspace{\intextsep}\hrule height .8pt\kern2pt
  \renewcommand{\caption}[2][]{\refstepcounter{algocf}\noindent\textbf{\algorithmcfname~\thealgocf:} ##2\par
    \kern2pt\hrule\kern2pt
    \ifx\relax##1\relax
      \addcontentsline{loa}{algocf}{\protect\numberline{\thealgocf}##2}\else
      \addcontentsline{loa}{algocf}{\protect\numberline{\thealgocf}##1}\fi
  }}{\par\kern2pt\hrule\addvspace{\intextsep}}
\makeatother
\usepackage{old-arrows}
\usepackage{color}
\usepackage{array, tabularx}
\usepackage{longtable}
\makeatletter
\@ifclassloaded{revtex4-2}{\appdef\@mkpream@relax{\let\@endpbox\relax}}{}
\makeatother
\usepackage{booktabs}
\newcolumntype{L}[1]{>{\raggedright\arraybackslash}p{#1}}
\newcolumntype{C}[1]{>{\centering\arraybackslash}p{#1}}
\newcolumntype{R}[1]{>{\raggedleft\arraybackslash}p{#1}}
\newcommand{\Circo}[0]{\mathfrak{C}}
\newcommand{\restrict}[2]{#1 \vert_{#2}}
\newcommand{\physq}[0]{\mathcal{M}}


\newcommand{\Xp}[0]{\mathcal{X}}
\newcommand{\Zp}[0]{\mathcal{Z}}

\newcommand{\nS}[1]{\mathbb{N}_{#1}}
\newcommand{\aS}[1]{\mathbb{C}_{#1}}

\newcommand{\nD}[1]{D^N_{#1}}
\newcommand{\aD}[1]{D^C_{#1}}

\newcommand{\control}[3]{C_{#1} {#2}_{#3}}
\usetikzlibrary{3d, shapes, arrows.meta, decorations.pathreplacing, decorations.markings, cd, shapes.geometric, calc}
\IfFileExists{tikzlibrarycalligraphy.code.tex}{\usetikzlibrary{calligraphy}}{}
\pgfkeysifdefined{/pgf/decoration/calligraphic brace/.@cmd}{}{\pgfkeys{/pgf/decoration/calligraphic brace/.style={brace}}}
\tikzstyle{midarrow} = [decoration={markings, mark= at position 0.6 with {\arrow{Stealth}}}, postaction={decorate}]
\renewcommand{\comment}[1]{{\color{red} #1}}

\newcommand{\YR}[1]{{\color{Teal} [Collin: #1]}}

\newcommand{\kc}[1]{{\color{cyan} [Katie: #1]}}

\newcommand{\bb}[1]{{\color{purple}[BB: #1]}}




\newcommand{\vva}[1]{{\color{magenta} [VVA: #1]}}









\newcommand{\al}[1]{{\color{DarkGreen} [Ana: #1]}}

\newcommand{\hp}[1]{{\color{brown} [Harry: #1]}}

\newcommand{\NT}[1]{{\color{orange} [NT: #1]}}

\newcommand{\dom}[1]{\textcolor{blue}{[Dom: #1]}}

\newcommand{\todo}[1]{\textcolor{blue}{[TO-DO: #1]}}

\newcommand{\new}[1]{{\color{DarkMagenta}#1}}
\definecolor{FGreen}{RGB}{58,158,92}
\definecolor{lightred}{RGB}{223, 123, 84}
\definecolor{lightbrown}{RGB}{242,216,199}
\definecolor{invisiblegray}{RGB}{238,238,237}
\definecolor{lightgreen}{RGB}{222, 235, 218}
\definecolor{lightcyan}{RGB}{153, 191, 216}
\definecolor{lightcyan2}{RGB}{208, 224, 235}
\definecolor{cyan2}{RGB}{144,201,226}
\definecolor{lightyellow}{RGB}{252, 238, 183}
\definecolor{OrangeRed}{RGB}{223,123,84}
\definecolor{CyanGreen}{RGB}{81,177,184}
\definecolor{rose}{RGB}{226,45,139}
\definecolor{Ycolor}{RGB}{248,236,104}
\definecolor{Hcolor}{RGB}{255,255,0}
\definecolor{Xcolor}{RGB}{246,143,138}
\definecolor{Zcolor}{RGB}{210,254,206}
\definecolor{Lred}{RGB}{244,163,163}
\definecolor{Lgreen}{RGB}{160,211,179}
\definecolor{Lblue}{RGB}{170,183,240}
\definecolor{tickcolor}{RGB}{181,0,250}
\definecolor{DarkBlue}{RGB}{14,38,80}
\definecolor{greenblue}{RGB}{0,255,127}
\definecolor{darkblue}{RGB}{57,85,134}
\definecolor{electrored}{RGB}{244,96,85}
\definecolor{electrored2}{RGB}{245,107,94}
\definecolor{BBXcheck}{RGB}{210,234,206}
\definecolor{BBZcheck}{RGB}{204,204,250}
\definecolor{softred}{RGB}{245,206,206}
\definecolor{fuchsiapurple}{RGB}{180, 0, 255}
\definecolor{Teal}{RGB}{62,138,141}
\definecolor{gold}{RGB}{233,168,58}
\definecolor{Crimson}{RGB}{220,20,60}
\definecolor{RedOrange}{RGB}{255,69,0}
\definecolor{DarkMagenta}{RGB}{139,0,139}
\definecolor{DarkGreen}{RGB}{0,100,0}
\definecolor{MidnightBlue}{RGB}{25,25,112}



\newcommand{\g}{\mathbf{g}}
\newcommand{\h}{\mathbf{h}}
\renewcommand{\k}{\mathbf{k}}
\newcommand{\f}{\mathbf{f}}
\newcommand{\C}{\mathbf{C}}
\newcommand{\Z}{\mathbb{Z}}
\renewcommand{\S}{\mathcal{S}}
\newcommand{\CC}{\mathcal{C}}
\newcommand{\charge}{k}
\newcommand{\totalcharge}{K}
\newcommand{\B}{\mathbb{B}}
\renewcommand{\O}{\mathcal{O}}
\newcommand{\Ocal}{\mathcal{O}}
\renewcommand{\Vec}{\mathcal{V}ec}
\newcommand{\Rep}{\text{Rep}}
\newcommand{\nnb}{\nonumber}
\newcommand{\lt}{\triangleleft}
\newcommand{\rt}{\triangleright}
\newcommand{\Lcal}{\mathcal{L}}
\newcommand{\Acal}{\mathcal{A}}
\newcommand{\Gcal}{\mathcal{G}}
\newcommand{\XX}{\mathcal{X}}
\newcommand{\YY}{\mathcal{Y}}
\newcommand{\ZZ}{\mathcal{Z}}
\newcommand{\Com}[2]{[ #1, #2 ]}
\newcommand{\ComGrp}[2]{[ #1, #2 ]_{\text{grp}}}
\newcommand{\level}{\operatorname{Lv}}
\newcommand{\code}{\mathcal{C}}


\newcommand{\glosym}[0]{\mathsf{L}}
\newcommand{\im}{\operatorname{im}}
\newcommand{\rk}{\operatorname{rk}}
\newcommand{\supp}[1]{\operatorname{Supp}(#1)}
\newcommand{\Ad}{\operatorname{Ad}}
\newcommand{\Irrep}{\operatorname{Irrep}}
\newcommand{\Wt}{\operatorname{Wt}}
\newcommand{\Lin}{\mathfrak{L}}
\newcommand{\interior}{\operatorname{int}}
\newcommand{\T}[0]{\mathsf{U}}
\newcommand{\orbit}[1]{O^{\mathcal{T}}(#1)}
\newcommand{\sym}[0]{\overline{S}}
\newcommand{\av}[0]{\mathbb{A}_v}
\newcommand{\bp}[0]{\mathbb{B}_f}
\newcommand{\onsite}[1]{\mathsf{T}|_{#1}}
\newcommand{\Log}[1]{\overline{#1}}
\newcommand{\id}[0]{\mathbf{1}}
\newcommand{\bd}[1]{\delta({#1})}
\newcommand{\zp}[0]{\mathbb{Z}_p}
\newcommand{\cX}[0]{\mathcal{X}}
\newcommand{\cZ}[0]{\mathcal{Z}}
\newcommand{\cC}[0]{\mathcal{C}}
\newcommand{\cE}[0]{\mathcal{E}}
\newcommand{\cx}[2]{CX_{#1\rightarrow#2}}
\newcommand{\wcp}[2]{\mathcal{C}P_{#1\rightarrow#2}}
\newcommand{\wcx}[2]{\mathcal{C}X_{#1\rightarrow#2}}
\newcommand{\wcz}[2]{\mathcal{C}Z_{#1\rightarrow#2}}
\newcommand{\wcy}[2]{\mathcal{C}Y_{#1\rightarrow#2}}
\newcommand{\cpw}[2]{C\mathcal{P}_{#1\rightarrow#2}}
\newcommand{\cxw}[2]{C\mathcal{X}_{#1\rightarrow#2}}
\newcommand{\czw}[2]{C\mathcal{Z}_{#1\rightarrow#2}}
\newcommand{\cyw}[2]{C\mathcal{Y}_{#1\rightarrow#2}}
\newcommand{\wct}[2]{\mathcal{C}\mathcal{T}_{#1\rightarrow#2}}
\newcommand{\wctdag}[2]{\mathcal{C}\mathcal{T}^\dagger_{#1\rightarrow#2}}
\newcommand{\cnot}[0]{\CNOT}
\newcommand{\SWAP}[0]{\mathsf{SWAP}}
\newcommand{\CNOT}[0]{\mathsf{CNOT}}
\newcommand{\CSWAP}[0]{\mathsf{CSWAP}}
\newcommand{\CX}[0]{\mathsf{CX}}
\newcommand{\CZ}[0]{\mathsf{CZ}}
\newcommand{\CT}[0]{\mathsf{CU}}
\newcommand{\swap}[0]{\SWAP}
\newcommand{\cswap}[0]{\CSWAP}

\newcommand{\overlap}[1]{\Omega(#1)}
\newcommand{\E}[2]{e_{(#1,#2)}}
\newcommand{\dec}[1]{\tilde{#1}}
\newcommand{\insuppof}[2]{#1|_{#2}}
\newcommand{\plushalf}[1]{{#1+\frac{1}{2}}}
\newcommand{\minushalf}[1]{{#1-\frac{1}{2}}}
\newcommand{\scrG}[0]{{\mathscr{G}}}
\newcommand{\sset}[1]{\Sigma(#1)}

\newcommand{\symOP}{\T}
\newcommand{\symS}{\overline{S}}
\newcommand{\haH}{\mathsf{H}}
\newcommand{\haA}{\mathsf{A}}
\newcommand{\haB}{\mathsf{B}}
\newcommand{\haP}{\mathsf{P}}
\newcommand{\haQ}{\mathsf{Q}}
\newcommand{\haX}{\mathsf{X}}
\newcommand{\Deco}{\mathfrak{D}}

\newcommand{\vertexset}[1]{\Sigma^{-1}(#1)}
\newcommand{\invvertexset}[1]{\Sigma(#1)}
\newcommand{\faceP}{\mathcal{P}} 
\newcommand{\vertexV}{\mathsf{V}}
\newcommand{\edgeE}{\mathsf{E}}
\newcommand{\faceF}{\mathsf{F}}
\newcommand{\graphG}{\mathsf{G}}
\newcommand{\intG}{\mathbb{G}}
\newcommand{\N}{\mathsf{N}}
\newcommand{\eigenE}{\mathbf{E}} 
\newcommand{\gaugedA}{\mathbb{A}}
\newcommand{\gaugedB}{\mathbb{B}}
\newcommand{\gaugedS}{\mathbb{S}}
\newcommand{\Sdec}{\mathfrak{s}}
\providecommand{\vr}{}
\renewcommand{\vr}{v_r}
\newcommand{\gaugedH}{\mathbb{H}}
\newcommand{\gaugingOP}{\mathscr{G}}
\newcommand{\chproj}[2]{\llbracket #1\rrbracket _{#2}}

\newcommand{\gX}{\mathcal{X}}
\newcommand{\gZ}{\mathcal{Z}}
\newcommand{\mX}{\mathcal{X}}
\newcommand{\mZ}{\mathcal{Z}}
\newcommand{\gW}{\mathcal{W}}
\newcommand{\gY}{\mathcal{Y}}
\newcommand{\Agrp}{\mathbf{A}}
\newcommand{\chOmega}{\mathit{\Omega}}
\newcommand{\Linex}{\Lin_{\rm ex}(\mathcal H_{\rm m})}
\newcommand{\Linnt}{\Lin_{\rm nt}^{\rm ex}(\mathcal H_{\rm mat})}


\newcommand{\K}{\mathcal{K}}
\newcommand{\gaugedT}{\mathbb{T}}
\newcommand{\Stab}{\mathcal{G}}
\newcommand{\F}{\mathcal{F}}
\newcommand{\Hp}{\mathcal{H}}
\newcommand{\Poly}{\mathsf{Poly}}
\newcommand{\Fphys}{F_{\text{phys}}}
\newcommand{\Fmeas}{F_{\text{meas}}}
\newcommand{\Fp}{\mathbf{F}_p}
\newcommand{\Zq}{\mathbb{Z}_q}
\newcommand{\Sprod}[2]{\Lambda\left( #1, #2 \right)}
\newcommand{\FT}{F^T}
\newcommand{\nbhd}{\mathsf{N}}

\newcommand{\Cgeom}{C_{\text{geom}}}
\newcommand{\CLDPC}{C_{\text{LDPC}}}
\newcommand{\Cstep}{C_{\text{step}}}
\newcommand{\Cretract}{C_{\text{retract}}}
\newcommand{\Ctotal}{C_{\text{total}}}
\newcommand{\Cgauge}{C_{\text{gauge}}}

\newcommand{\nonabel}[1]{S^\mathrm{N}_{#1}}
\newcommand{\abel}[1]{S^\mathrm{A}_{#1}}
\newcommand{\gnonabel}[1]{\mathbb{S}^\mathrm{N}_{#1}}
\newcommand{\gabel}[1]{\mathbb{S}^\mathrm{A}_{#1}}
\newcommand{\onset}[2]{{#1}_{\sset{#2}}}
\newcommand{\ccomm}[2]{[#1,#2]_{\mathrm{grp}}}
\newcommand{\tu}[0]{t_{\mathrm{u}}}
\newcommand{\tg}[0]{t_{\mathrm{g}}}
\newcommand{\tf}[0]{t_\mathrm{f}}
\newcommand{\Ig}[0]{\mathrm{I}_{\mathrm{g}}}
\newcommand{\Ifn}[0]{\mathrm{I}_\mathrm{f}}
\newcommand{\Ib}[0]{\mathrm{I}_0}
\newcommand{\cA}[0]{\mathcal{A}}
\newcommand{\cQ}[0]{\mathcal{Q}}
\newcommand{\cK}[0]{\mathcal{K}}
\newcommand{\grpcom}[2]{\ComGrp{#1}{#2}}
\newcommand{\de}[0]{d_\mathcal{E}}



\definecolor{Teal}{RGB}{0,128,128}
\definecolor{Crimson}{RGB}{220,20,60}
\definecolor{RedOrange}{RGB}{255,69,0}
\definecolor{DarkMagenta}{RGB}{139,0,139}
\definecolor{DarkGreen}{RGB}{0,100,0}
\definecolor{MidnightBlue}{RGB}{25,25,112}


\hypersetup{
    colorlinks=true,
    linkcolor=MidnightBlue,
    citecolor=MidnightBlue,
    urlcolor=MidnightBlue,}

\usepackage{amsfonts,latexsym,amsthm,amssymb,amsmath,amscd,euscript,tikz,mathtools}
\usepackage{framed}
\usepackage[normalem]{ulem}
\makeatletter
\@ifclassloaded{revtex4-2}{}{\usepackage[margin=1in]{geometry}}
\makeatother
\usepackage{color}




\usepackage{enumitem}
\usepackage{environ}
\usepackage{dcolumn}

\usetikzlibrary{positioning, shadows, arrows, backgrounds, fit}
\usetikzlibrary{decorations.pathmorphing}

\definecolor{Teal}{RGB}{0,128,128}

\hypersetup{
    colorlinks=true,
    linkcolor=MidnightBlue,
    citecolor=MidnightBlue,
    urlcolor=MidnightBlue,
}

\definecolor{GrayOut}{gray}{0.8}
\NewEnviron{GrayedOut}{\begingroup\color{GrayOut}\BODY\endgroup
}

\newcommand{\diam}{\operatorname{diam}}
\newcommand{\dist}{\operatorname{dist}}

\newcommand{\cH}{\mathcal{H}}
\newcommand{\cT}{\mathcal{T}}
\newcommand{\cN}{\mathcal{N}}
\newcommand{\cB}{\mathcal{B}}
\newcommand{\bbP}{\mathbb{P}}
\newcommand{\haT}{\mathsf{T}}
\newcommand{\haD}{\mathsf{D}}
\newcommand{\cF}{\mathcal{F}}

\newcommand{\CZero}{C_0(\Gcal_0)}
\newcommand{\COne}{C_1(\Gcal_1)}
\newcommand{\deltaOne}{\delta_1}
\newcommand{\jall}{\mathbf{1}}
\newcommand{\tof}[0]{\ket{\mathsf{Tof}}}
\newcommand{\Pp}[1]{\mathrm{P}^+_{\Log{#1}}}
\newcommand{\Ppk}[1]{\mathrm{P}^+_{\Log{#1}^{\otimes k}}}
\newcommand{\Pm}[1]{\mathrm{P}^-_{\Log{#1}}}

\newcommand{\fa}[0]{F_\mathrm{A}}
\newcommand{\fn}[0]{F_\mathrm{N}}

\newcommand{\CHANGE}[1]{{\color{DarkMagenta}#1}}

\newcommand{\syndromeSet}{\sigma}
\newtheorem{axiom}{Axiom}[section]
\newtheorem{question}{Question}[section]


\newcommand{\green}[1]{\textcolor{Teal}{#1}}
\newcommand{\blue}[1]{\textcolor{blue}{#1}}
\newcommand{\red}[1]{\textcolor{red}{#1}}

\makeatletter
\let\UncappedLeft\left
\let\UncappedRight\right
\def\left{\@ifnextchar({\CappedLeftParen}{\UncappedLeft}}
\def\CappedLeftParen(#1\right{\@ifnextchar){\CappedRightParen{#1}}{\UncappedLeft(#1\UncappedRight}}
\def\CappedRightParen#1){\mathopen{}\bigg(#1\bigg)\mathclose{}}
\makeatother

\definecolor{vred}{rgb}{0.85,0,0}
\definecolor{epur}{cmyk}{0.45,0.86,0,0}
\definecolor{poraBase}{cmyk}{0,0.42,1,0}
\colorlet{pora}{poraBase!90!black}
\definecolor{setgray}{RGB}{225,225,225}
\newcommand{\vC}[1]{{\color{vred}#1}}
\newcommand{\eC}[1]{{\color{epur}#1}}
\newcommand{\pC}[1]{{\color{pora}#1}}
\newcommand{\sq}[9]{\left(\begin{smallmatrix}
\vC{#1} & \eC{#2} & \vC{#3}\\[1pt]
\eC{#4} & #5 & \eC{#6}\\[1pt]
\vC{#7} & \eC{#8} & \vC{#9}
\end{smallmatrix}\right)}
\newcommand{\mm}[4]{\begin{smallmatrix}#1 & #2\\ #3 & #4\end{smallmatrix}}
\newcommand{\sqe}[5]{\sq{\cdot}{#1}{\cdot}{#2}{#3}{#4}{\cdot}{#5}{\cdot}}
\newcommand{\Sg}{\mathsf{S}}
\newcommand{\CZbar}{\overline{CZ}}

\setcounter{tocdepth}{1}

    
    \author{}
    \title{Algebraic formalism and distance of the gauged code}
    
    \begin{document}
    
    
    \maketitle

\fi

The gauging measurement protocol was defined in the section above, it remains to prove that this protocol is \emph{fault-tolerant}. Here, ``fault-tolerance'' means that the minimum un-detectable fault that causes a logical failure is of large weight. This includes the application of a physical logical operator to the code, and inferring an incorrect measurement outcome of the logical Clifford $\T$. As an intermediate step to the proof of large fault distance in both space and time, we now focus on characterizing the properties of the gauged code. Specifically, we demonstrate that the gauged code's distance satisfies $d_g\geq d/\nu$, where $\nu$ is the maximum size of a vertex set, so long as the graph desiderata in Def.~\ref{def:graph-desiderata} are satisfied. 

To enable the proof, we introduce an alternate, algebraic formulation of the gauging procedure from Section \ref{sec:gauging-procedure}, following Refs.~\cite{Bulmash_2019,Prem_2019}. After re-casting the gauging protocol in this language, we demonstrate that the gauged codespace and the $\T$-symmetric subspace of the original stabilizer code are isomorphic. This isomorphism facilitates the proof of the gauged code distance in Theorem \ref{thm:space-distance}.
 
\begin{table}[H]
    \centering
    \begin{tabular}{@{}p{0.30\linewidth}p{0.67\linewidth}@{}}
    \hline
        Notation & Denotes \\
    \hline
        $\Agrp = \langle \T_1, \dots \T_n \rangle$ & Finite Abelian group generated by symmetries being gauged. \\
        $\widehat{\Agrp}$ & Irreducible representations of $\Agrp$.\\
        $\graphG = (\vertexV, \edgeE, \faceF)$ & Gauging graph and associated vertices, edges, and faces. \\
        $\Agrp[\vertexV], \Agrp[\edgeE], \Agrp[\faceF]$ & Assignments of elements of $\Agrp$ to every vertex, edge, or face. \\
        $\widehat{\Agrp}[\vertexV], \widehat{\Agrp}[\edgeE], \widehat{\Agrp}[\faceF]$ & Assignments of elements of $\widehat{\Agrp}$ to every vertex, edge, or face. \\
        $\delta_i, \partial_i$ & Coboundary and boundary maps \\
        $\T^g$ & Clifford symmetry corresponding to group element $g \in \Agrp$.\\
        $\T^{\bm{a}}$ & Product of local factors $\T_v$ according to the assignment $\bm{a} \in \Agrp[\vertexV]$\\
        $\Xp, \Zp$ & Generalized Pauli operators for $\Agrp$\\
        $\Xp^{\bm{a}}$ & Product of $\Xp$ according to the group element assignment $\bm{a}$.\\
        $\Zp^{\bm{\chi}}$ & Product of $\Zp$ according to the irrep assignment $\bm{\chi}$. \\
        $\mathcal{H}_{\rm{sym.}}$ & Subspace determined by $\T = +1$. \\
        $\mathcal{H}_{\rm{gau.}}$ & Subspace determined by $\av = +1$, $\bp = +1$ for all vertices and faces. \\
        $\Gamma, \tilde{\Gamma}$ & The gauging map introduced in this section and the gauging circuit (Section \ref{sec:gauging-procedure}). \\
        $h_{\Agrp}(\graphG)$ & Generalized Cheeger constant of $\graphG$.
        
    \end{tabular}
    \caption{Notation used for this section.}
    \label{tab:sec-3-notation}
\end{table}

\subsection{Algebraic formalism for gauging}

We begin by reformulating the gauging procedure discussed in Sec.~\ref{sec:gauging-procedure} in terms of algebraic operations. Suppose we are gauging some set of commuting Clifford symmetries $\{\T_i\}$ with $\T_i^{p_i} = \id$. These Clifford operators generate a finite Abelian group $\Agrp = \times_i \mathbb{Z}_{p_i}$. For example, for symmetries satisfying $\T^2 = \id$, the corresponding Abelian group is $\Agrp = \mathbb{Z}_2$.

\begin{definition}[Gauging 2-complex]
\label{def:gauging-complex}
Let $\vertexV$, $\edgeE$, $\faceF$ denote the vertices, edges and faces of the gauging graph $\graphG=(\mathsf{V},\mathsf{E})$ from Section~\ref{sec:gauging-procedure}. For any finite Abelian group $A$, let $\Agrp[\vertexV]$, $\Agrp[\edgeE]$, $\Agrp[\faceF]$ denote assignments of an element of $A$ to every vertex, edge and face,  respectively.
The relevant coboundary maps are 
\begin{align}
    \graphG^\bullet = \Agrp[\vertexV] \mathop{\rightarrow}^{\delta_{1}} 
    \Agrp[\edgeE] \mathop{\rightarrow}^{\delta_2} 
    \Agrp[\faceF] .
\end{align}
Similarly, let $\widehat{\Agrp}[\vertexV]$, $\widehat{\Agrp}[\edgeE]$, $\widehat{\Agrp}[\faceF]$ denote assignments of an irreducible representation label of $\Agrp$ to every vertex, edge and face, respectively.
The relevant boundary maps are
\begin{align}
    \graphG_\bullet = \widehat{\Agrp}[\vertexV] \mathop{\leftarrow}^{\partial_{1}} 
    \widehat{\Agrp}[\edgeE] \mathop{\leftarrow}^{\partial_2} 
    \widehat{\Agrp}[\faceF] .
\end{align}
\end{definition}


We use $+$ to denote the group operation in $\Agrp$, and bold symbols to denote assignments of group element to vertices, edges or faces.
If $\bm{a}, \bm{b}\in \Agrp[\vertexV]$, then their sum, $\bm a+\bm b $ is defined vertex-wise $(\bm a + \bm b)_v =  a_v + b_v$.
For a group element $g \in \Agrp$, we abuse notation and write $ g \id \in \Agrp[\vertexV]$ for the constant vertex assignment where $(g\id)_v = g$ for every $v\in \vertexV$. We write $g \bm{e_v} \in \Agrp[\vertexV]$ for the assignment with group element $g$ on vertex $v \in \vertexV$, and the identity element $0$ elsewhere. Similarly $g\bm{e_f} \in \Agrp[\faceF]$ is the assignment with $g$ on face $f \in \faceF$ and identity elsewhere.

Furthermore, we define $\T^g$ where $g \in \Agrp$ as the logical Clifford operator representing the group element $g$. In the case $\Agrp = \mathbb{Z}_q = \langle \T ~|~ \T^p = \id \rangle$ for $p\geq 2$ is a cyclic group, we have $g = 0, \dots, p-1$. Then $\T^g$ is simply the sole generator $\T$ multiplied $g$ times. In general, if we have Abelian group $\Agrp = \mathbb{Z}_{p_1} \times \cdots \times \mathbb{Z}_{p_n}$ we define:
\begin{equation}
    \T^g = \prod_{i=1}^n \T_i^{g_i}, \quad g = (g_1, \dots, g_n) \in \Agrp.
\end{equation}
When we write $\T^{\bm{a}}$, where the superscript is bolded, this denotes the multiplication of different local symmetry operators for each vertex set, determined by the value of $\bm{a} \in \Agrp[\vertexV]$ at each vertex.
\begin{equation}
    \T^{\bm{a}} = \prod_{v \in \vertexV} (\T_v)^{a_v}, \quad a_v =((a_v)_1, \dots, (a_v)_n) \in \Agrp
\end{equation}

We also generalize the clock and shift operators defined for $\Agrp = \mathbb{Z}_p$ to the case $\Agrp = \mathbb{Z}_{q_1} \times \cdots \times \mathbb{Z}_{q_n}$. For each cyclic group component $\mathbb{Z}_{p_i}$, the shift operator is $\Xp^{(i)}$ and the clock operator is $\Zp^{(i)}$. Shift operators are associated with group elements (through the $+$ group operation) while clock operators evaluate an irreducible representation of $\mathbb{Z}_{p_i}$ on a given state.
\begin{equation}
    \Zp^{(i)} \Xp^{(i)} = \omega_i \Xp^{(i)} \Zp^{(i)}, \quad \omega_i = \exp{2\pi i/p_i}
\end{equation}
The clock and shift operators for the full group are constructed from the tensor product of these elementary generators, with the corresponding commutation relation:
\begin{equation}
    \Xp \Zp = \omega \Zp \Xp, \quad \omega = \prod_{i = 1}^n \omega_i = \exp{2 \pi i \sum_{i = 1}^n 1/p_i} .
\end{equation}

The generalized Pauli operators $\Xp$ and $\Zp$ can be labeled by bolded superscripts indicating their support on the gauging graph. These superscripts may be elements of $\Agrp[\vertexV]$ or $\Agrp[\edgeE]$, if we are discussing Pauli operator support on the vertex ancillas or the edge ancillas respectively. For example, a generic configuration of shift operators on the vertex ancilla may be written:
\begin{equation}
    \Xp^{\bm{a}} = \prod_{v \in \vertexV} \Xp_v^{a_v}, \quad \bm{a} \in \Agrp[\vertexV]
\end{equation}
while a set of clock operators supported on the edges is labeled by a choice of irreducible representation for each edge:
\begin{equation}
    \Zp^{\bm{b}} = \prod_{e \in \edgeE} \Zp_e^{b_e}, \quad \bm{b} \in \widehat{\Agrp}[\edgeE]
\end{equation}

We call the symmetric subspace of our initial stabilizer code $\cC$ under $\T$ the \textit{symmetric code}. 

\begin{definition}[Symmetric code]
\label{def:symmetric-code}
Let
\begin{align}
    \mathcal H_{\rm sym.} := \{\ket{\psi} : \T^a\ket{\psi}=\ket{\psi} \quad \forall a\in \Agrp\}.
\end{align}
be the symmetric subspace. The projector
\begin{align}
    \Pi_{\rm sym.} = \frac{1}{|\Agrp|} \sum_{a \in \Agrp} \T^a
\end{align}
projects onto $\mathcal H_{\rm sym.}$. The symmetric code is the intersection of this subspace with the codespace of the initial code, $\cC$. The symmetric code is defined by $\T$ and the symmetrized versions of the original code stabilizers: 
\begin{equation}
    \sym = \Pi_{\rm sym.} S \Pi_{\rm sym.} = \frac{1}{|\Agrp|} \sum_{g \in \Agrp} (\T^g)^\dagger S \T^g.
\end{equation}
\end{definition}

As we demonstrate below, the gauging protocol, when viewed as a formal mathematical object, is an isomorphism between the symmetric code and the gauged code $\cC'$. A convenient description of the gauging map for demonstrating this isomorphism is the following:

\begin{definition}[Gauging map]
\label{def:gauging-map}
Let $\mathcal{H}_\mathrm{m}$ be the Hilbert space of the data qudits of the initial code $\cC$, and let $\mathcal{H}_\mathrm{g}$ be the Hilbert space of gauge qudits on edges $e\in \edgeE$.
For $\ket{\psi}\in \mathcal H_{\rm sym.}$, we define the gauging map $\Gamma:\mathcal{H}_\mathrm{sym}\rightarrow \mathcal{H}_\mathrm{m}\otimes \mathcal{H}_\mathrm{g}$. 
\begin{align}
\label{eq:gauging-map}
     \Gamma \ket{\psi} &:= \frac{1}{|\Agrp|^{(|\vertexV|-1)/2}}\sum_{\bm{a}\in \Agrp[\vertexV]/\Agrp} \T^{\bm{a}} \ket{\psi}\ket{\delta_1\bm{a}}
\end{align}
\end{definition}

The derivation of equation~\eqref{eq:gauging-map} is the following.
Start with the state $\ket{\psi}\ket{0}^{\otimes|\edgeE|}$, where every edge qudit is initialized in the $\ket{0}$ state.
For each vertex $v\in \vertexV$, the projector onto the $\av=1$ subspace is 
\[
\Pi_{\av}^{+} = \frac{1}{|\Agrp|} \sum_{g\in \Agrp} \av^g, \quad \av^g = \T_v^g \otimes \Xp^{\delta_1 g \bm{e_v}}
\]
Multiplying these projectors over all vertices $v$, we have
\[
\prod_{v\in \vertexV} \Pi_{\av}^{+} = \prod_{v\in \vertexV} \frac{1}{|\Agrp|} \sum_{g\in \Agrp} \av^g= \frac{1}{|\Agrp|^{|\vertexV|}}\sum_{\bm{a}\in \Agrp[\vertexV]} A(\bm{a})
\]
where $\bm{a}=(a_v)_{v\in \vertexV}\in \Agrp[\vertexV]$ runs over every possible choice of assigning a group element to each vertex, and $A(\bm{a}):=\prod_{v\in \vertexV} \av^{a_v}=\T^{\bm{a}}\otimes \Xp^{\delta_1\bm{a}}$.
It follows that 
\begin{align}
    \prod_{v\in \vertexV} \Pi_{\av}^{+}\ket{\psi}\ket{0}^{\otimes |\edgeE|} = \frac{1}{|\Agrp|^{|\vertexV|}}\sum_{\bm{a}\in \Agrp[\vertexV]}
    \T^{\bm{a}}\ket{\psi}\ket{\delta_1\bm{a}}.
    \label{eq:unsimplified-gauging}
\end{align}
A further simplification can be performed when multiple $\bm{a}$ have the same coboundary ($\delta_1\bm{a}=\delta_{1}\bm{a}'$).
In particular, this is true for $\bm{a}'=\bm{a}+g\id \ \forall g\in \Agrp$ because $g\id\in \Agrp[\vertexV]$ has the same value from both ends of each edge in a fully connected graph, $\delta_1(g\id)=0$.
In other words, two vertex labelings have the same coboundary when they differ by a constant $g\id$.
Since the states are identical when two labelings have the same coboundary, the final, normalized state after applying the gauging map is Eq.~\eqref{eq:gauging-map}.

Let $\bm{\bm{\chi}}=(\bm{\chi}_v)_{v\in \vertexV}\in \widehat{\Agrp}[\vertexV]$ be an assignment of an irreducible representation of $\Agrp$ to each vertex.
For $\bm{a}=(a_v)_{v\in\vertexV}\in \Agrp[\vertexV]$, define 
\begin{align}
    \bm{\bm{\chi}}(\bm{a}) :=\prod_{v\in \vertexV} \bm{\chi}_v(a_v),
\end{align}
which is the product of vertexwise eigenvalues associated with group element configuration $\bm{a}$.

We now define the gauged form of a matter operator $M$, meaning an operator supported entirely on matter qudits. We first decompose any $M$ into terms $M_{\bm{\bm{\chi}}}$, each of which has eigenvalue $\bm{\bm{\chi}}(\bm{a})$ under conjugation by $\T^{\bm{a}}$ for every $\bm{a}\in \Agrp[\vertexV]$\footnote{For readers familiar with the condensed matter theory literature, $\bm{\bm{\chi}}(\bm{a})$ is the charge configuration of the operator $M$ under the action of $\T^{\bm{a}}$.}.
We refer to this decomposition as the eigenoperator decomposition of a matter operator $M$.


\begin{definition}[Eigenoperator decomposition] \label{def:eigen-decomp}

For any matter operator $M$, its \emph{eigenoperator components} $M_{\bm{\bm{\chi}}}$ are
\begin{align}
    M_{\bm{\bm{\chi}}} := \frac{1}{|\Agrp|^{|\vertexV|}}\sum_{\bm{a}\in\Agrp[\vertexV]}\bm{\bm{\chi}}(\bm a)^*\T^{\bm{a}}M (\T^{\bm{a}})^\dagger,
    \label{eq:eigencomponent}
\end{align}
where each eigenoperator component satisfies
\begin{align}
    \T^{\bm{a}}M_{\bm{\bm{\chi}}} (\T^{\bm{a}})^\dagger = \bm{\bm{\chi}}(\bm a ) M_{\bm{\bm{\chi}}}.
\end{align}
Then the \emph{eigenoperator decomposition} of $M$ is
\begin{align}
    M=\sum_{\bm{\bm{\chi}} }M_{\bm{\bm{\chi}}}.
    \label{eq:eigenoperator-decomposition}
\end{align}
\end{definition}
One can verify that the $M_{\bm{\bm{\chi}}}$ defined in Eq.~\eqref{eq:eigencomponent} are orthogonal components of $M$ and satisfy $\sum_{\bm{\bm{\chi}}} M_{\bm{\bm{\chi}}}=M$. Once an operator is decomposed into its eigencomponents, it is straightforward to apply the gauging map.

\begin{definition}[Gauged Operator]
\label{def:gauged-operator}
Let $M=\sum_{\bm{\bm{\chi}}}M_{\bm{\bm{\chi}}}$ be a matter operator such that 
$M_{\bm{\bm{\chi}}} = 0$ when $\bm{\bm{\chi}}\notin\mathop{Im} (\partial_1)$. In other words,
\begin{align}
M=\sum_{\bm{\bm{\chi}} \in \mathop{Im} \partial_1} M_{\bm{\bm{\chi}}}  .
\end{align}
For each $\bm{\bm{\chi}} \in \mathop{Im} (\partial_1)$, choose a $\bm{p}_{\bm{\bm{\chi}}}\in \widehat{\Agrp}[\edgeE]$, or a path of edges labeled by irreps of $A$ such that
\begin{align}
    \partial_1\bm{p}_{\bm{\bm{\chi}}} = \bm{\bm{\chi}}.
\end{align}

We define the \emph{gauged operator} of $M$ as $\Gamma[M]$:
\begin{align}
    \Gamma[M]&= \sum_{\bm{\bm{\chi}} \in \mathop{Im} \partial_1} M_{\bm{\bm{\chi}}} \otimes \Zp^{\bm{p}_{\bm{\bm{\chi}}}}
\end{align}


\end{definition}

We have fully introduced the algebraic gauging formalism; however, we have not yet demonstrated the equivalence with the gauging circuit we introduced in Section \ref{sec:gauging-procedure}. 

\begin{lemma}[Equivalence of the gauging circuit and gauging map] \label{lem:gauging-circuit-map-equivalence}
    Consider the gauging circuit $\widetilde{\Gamma}$ introduced in the previous section (Definition \ref{def:gauging-circuit}). This circuit acts identically to the gauging map defined in Def. \ref{def:gauged-operator} within the $\bp = +1$ subspace. 
\end{lemma}
\begin{proof}
    Consider the action of the gauging circuit on a single eigen-operator component $M_{\bm{\bm{\chi}}}$.
    \begin{equation}
        \begin{aligned}
        \widetilde{\Gamma} = \underbrace{\left (\prod_{v \in \vertexV} \dyad{+}_v \right )}_\textrm{P} \underbrace{\left ( \prod_{v \in \vertexV} \prod_{e \in \delta v} C_v \Xp^{\eta_v(e)}_e\right )}_{\textrm{C.E.}} \underbrace{\left ( \prod_{v \in \vertexV} C_v \T_v \right )}_\textrm{S.E.} \underbrace{\left ( \prod_{v \in \vertexV} \ket{+}_v \prod_{e \in \edgeE} \ket{0}_e \right )}_\textrm{Ancilla Init.} 
        \end{aligned}
    \end{equation}
    \begin{equation}
        \begin{aligned}
        \widetilde{\Gamma} M_{\bm{\bm{\chi}}} \widetilde{\Gamma}^\dagger &= \mathrm{P} (\mathrm{C.E.}) (\mathrm{S.E.}) (\mathrm{Init.}) M_{\bm{\bm{\chi}}} (\mathrm{Init.})^\dagger(\mathrm{S.E.})^\dagger   (\mathrm{C.E.})^\dagger \mathrm{P}^\dagger \\
        &= \mathrm{P} (\mathrm{C.E.}) (\mathrm{S.E.}) M_{\bm{\bm{\chi}}} (\mathrm{S.E.})^\dagger   (\mathrm{C.E.})^\dagger \mathrm{P}^\dagger \\
        &= \mathrm{P} (\mathrm{C.E.}) [ M_{\bm{\bm{\chi}}} \otimes \Zp(\bm{\bm{\chi}})] (\mathrm{C.E.})^\dagger \mathrm{P}^\dagger \\
        &= \mathrm{P}  [ M_{\bm{\bm{\chi}}} \otimes \Zp^{\bm{\bm{\chi}}} ] \mathrm{P}^\dagger \\
        &= M_{\bm{\bm{\chi}}} \otimes \Zp^{\bm{p}_{\bm{\bm{\chi}}}}, \quad \partial_1 \bm{p}_{\bm{\bm{\chi}}} = \bm{\bm{\chi}} 
        \end{aligned}
    \end{equation}
    In the last step, we have used the edge stabilizers generated during the cluster entangling step to push $\Zp^{\bm{\bm{\chi}}}$ off the vertices onto the edges. This leads to the constraint $\partial_1 \bm{p}_{\bm{\bm{\chi}}} = \bm{\bm{\chi}}$. By linearity, we must have 
    \begin{equation}
        \widetilde{\Gamma}[M] = \sum_{\bm{\bm{\chi}} \in \mathrm{Im}\partial_1} \widetilde{\Gamma}[M_{\bm{\bm{\chi}}}] = \sum_{\bm{\bm{\chi}} \in \mathrm{Im}\partial_1} M_{\bm{\bm{\chi}}} \otimes  \Zp^{\bm{p}_{\bm{\bm{\chi}}}} \stackrel{\bp = +1}{=} \Gamma[M].
    \end{equation}
    In the last equality, emphasize that the action is equivalent inside the $\bp = +1$ space, as we may choose different path conventions.
\end{proof}

Recall the form of the gauged code obtained via the gauging circuit; Lemma \ref{lem:gauging-circuit-map-equivalence} assures that one obtains an equivalent code using the gauging map.
\begin{definition}[Gauged code]
\label{def:gauged-code-stabs}
    For ease of the reader, we copy the gauged code stabilizers derived in Section \ref{sec:gauging-procedure} here. We have included a translation of the stabilizers into the notation introduced for generic Abelian groups $\Agrp$.

    \begin{equation}
        \av^g = \T_v^g \Xp^{\delta_1 g \bm{e_v}}, \quad \bp^g = \Zp^{\partial_2 g \bm{e_f}}, \quad \bar{\mathbb{S}}  = \Gamma[\sym] 
    \end{equation}
    
    \begin{equation}
        \Gamma[\sym] = \sum_{\bm{\bm{\chi}} \in \mathop{Im} \partial_1} \sym_{\bm{\bm{\chi}}} \otimes \Zp^{\bm{p}_{\bm{\bm{\chi}}}}
    \end{equation}
    where $\partial_1\bm{p}_{\bm{\bm{\chi}}} = \bm{\bm{\chi}}.$
\end{definition}

\begin{remark}[Path independence of the gauging map]

We note that while the explicit form of the gauged operators presented in Definition \ref{def:gauged-operator} is dependent on a choice of edges $\bm{p}_{\bm{\bm{\chi}}}$, any other valid choice $\bm{q}_{\bm{\bm{\chi}}}$ differs by a cycle: 
\begin{equation}
    \partial_1 (\bm{p}_{\bm{\bm{\chi}}} - \bm{q}_{\bm{\bm{\chi}}}) = \bm{\bm{\chi}} - \bm{\bm{\chi}} = 0 
\end{equation}
Recall this cycle $\partial_2 \bm{f} = \bm{p}_{\bm{\bm{\chi}}} - \bm{q}_{\bm{\bm{\chi}}} \in \widehat{\Agrp}[\edgeE]$ is associated to a cycle stabilizer of the gauged code $\mathbb{B}_f = \Zp^{\partial_2 \bm{f}}$. Thus, in the subspace where all cycle stabilizers $\mathbb{B}_f = +1$, all paths with the same boundary are equivalent choices. In practice, there are other considerations: for example, when writing down the decorated stabilizers, we want our path choices to be constant weight for a given stabilizer generator, to ensure LDPC-ness of the gauged code.
\end{remark}


\subsection{The gauging map is an isomorphism}

We now prove an isomorphism between the $+1$ eigenspace of $\T$, $\mathcal{H}_{\rm{sym.}}$, and the subspace determined by $\av = \bp = +1$ constraints, $\mathcal{H}_{\rm{gau.}}$. The proof follows from basic algebraic manipulations, using the form of the gauging map in Def. \ref{def:gauging-map}. 


\begin{lemma}[The gauging map is an isomorphism]
\label{lem:gauging-map-isomorphism}
    Let $\graphG$ be connected so $H^1(\graphG^\bullet)=0$. 
    Then, $\Gamma$ is an isomorphism,
    \begin{align}
        \Gamma:\mathcal H_{\rm sym.}\xrightarrow[]{\cong} \ \mathcal{H}_{\rm gau.} \equiv \{\ket{\phi}: \av\ket{\phi}=\ket{\phi}\ \forall v,\; \bp\ket{\phi}=\ket{\phi}\ \forall p\}, 
    \end{align}
    and 
    \begin{align}
        \Gamma[M]\Gamma\ket{\psi} = \Gamma(M\ket{\psi})
        \label{eq:gauging-map-associativity}
    \end{align}
    for every $M$ for which $\Gamma[M]$ in Def.~\ref{def:gauged-operator} is defined and every $\ket{\psi}\in \mathcal H_{\rm sym.}$ (Def.~\ref{def:symmetric-code}).
\end{lemma}
\begin{proof}
    First we have that 
\begin{align}
     \Gamma^\dagger \Gamma &= \frac{1}{|\Agrp|^{|\vertexV|+1}} \sum_{\bm{a}, \bm{b} \in \Agrp[\vertexV]} (\T^{\bm{b}})^\dagger \T^{\bm{a}} \otimes \langle \delta_1 \bm{b} \ket{\delta_1 \bm{a}} \\
     &= \frac{1}{|\Agrp|^{|\vertexV|+1}} \sum_{\bm{a} \in \Agrp[\vertexV]} \sum_{\substack{\bm{b} \text{ s.t. } \\ \delta_1 (\bm{b} - \bm{a}) = 0}} \T^{\bm{a} - \bm{b}} \label{eq:isomorphism-1}\\
     &= \frac{1}{|\Agrp|^{|\vertexV|+1}} \sum_{\bm{a} \in \Agrp[\vertexV]} \sum_{g \in \Agrp} \T^g\label{eq:isomorphism-2}\\
     &= \frac{1}{|\Agrp|} \sum_{g \in \Agrp} \T^g = \Pi_{\rm sym.}\label{eq:isomorphism-3}
\end{align}
    From Eq.~\eqref{eq:isomorphism-1} to Eq.~\eqref{eq:isomorphism-2}, we have used that $\graphG$ is connected and therefore $\delta_1\bm a =\delta_1\bm b$ iff $\bm a =\bm b + c\id$ for some $c\in \Agrp$. Eq.~\eqref{eq:isomorphism-3} means that acting the gauging map and then undoing it implements a projection onto the symmetric subspace $\mathcal H_{\rm sym.}$. Furthermore, it follows that the gauging map preserves the inner product between any two $\ket{\psi},\ket{\phi}\in \mathcal H_{\rm sym.}$, since $\bra{\phi}\Gamma^\dagger\Gamma\ket{\psi} = \bra{\phi}\Pi_{\rm sym.}\ket{\psi}=\langle\phi\ket{\psi}$.

    The reverse procedure: starting from the gauged degrees of freedom, undoing the gauging, and then re-doing it, yields the projector onto the subspace with $\av = +1$ and $\bp = +1$: 
\begin{equation}
\begin{aligned}
    \Gamma \Gamma^\dagger
    &= \frac{1}{|\Agrp|^{|\vertexV|+1}} \sum_{\bm{a}, \bm{b} \in \Agrp[\vertexV]} \T^{\bm{a}-\bm{b}}  \otimes \ket{\delta_1 \bm{a}} \bra{\delta_1 \bm{b}} \\
    &=  \frac{1}{|\Agrp|^{|\vertexV|+1}} \sum_{\bm{a}, \bm{b} \in \Agrp[\vertexV]} \T^{\bm{a}-\bm{b}}  \otimes \Xp^{\delta_1 (\bm{b} - \bm{a})} \dyad{\delta_1 \bm{b}} \\
    &= \frac{1}{|\Agrp|^{|\vertexV|}} \left ( \sum_{\bm{a'}} (\T^{\bm{a'}})^\dagger \otimes \Xp^{\delta_1\bm{a'}}\right ) \frac{1}{|\Agrp|} \left ( \sum_{\bm{b}} \dyad{\delta_1 \bm{b}} \right ) \\
    &= \prod_{v\in \vertexV} \Pi_{\av}^{+} \prod_{p } \Pi_{\bp}^+
\end{aligned}
\end{equation}
Here we have defined the projector onto all $\bp = +1$ as follows:
    \begin{equation}
    \begin{aligned}
        \prod_{f \in \faceF} \Pi_{\bp}^+ &= \prod_{f \in \faceF} \frac{1}{|\Agrp|} \sum_{\chi \in \widehat{A}} \Zp^{\partial_2 \bm{e_f}} \\
        &= \frac{1}{|\Agrp|^{|\faceF|}} \sum_{\bm{\chi} \in \widehat{A}[\faceF]} \Zp^{\partial_2 \bm{\chi}}\\
        &= \sum_{\bm{\epsilon} \in \Agrp[\edgeE]}\left( \frac{1}{|\Agrp|^{|\faceF|}}  \sum_{\bm{\chi} \in \widehat{A}[\faceF]} \omega^{\bm{\chi}(\bm{\epsilon})} \right )\dyad{\bm{\epsilon}}\\
        &=\sum_{\bm{\epsilon} \in \Agrp[\edgeE]} [0 \text{ unless } \bm{\epsilon} \in \text{Im} \delta_1]\dyad{\bm{\epsilon}}\\
        &=\frac{1}{|\Agrp|} \sum_{\bm{a} \in \Agrp[\vertexV]} \dyad{\delta_1 \bm{a}}
    \end{aligned}
\end{equation}
In the last line, we have divided by $|\Agrp|$ to remove over-counting of configurations $\bm{a}$ that differ by an element of $\ker \delta_1$. Thus $\Gamma$ is an isomorphism between the symmetric subspace and the subspace generated by the $\av$ and $\bp$ stabilizers. 

Now we show Eq.~\eqref{eq:gauging-map-associativity}.
We compute
\begin{align}
    \Gamma[M]\Gamma\ket{\psi}&= \frac{1}{|\Agrp|^{(|\vertexV|-1)/2}}\sum_{\bm a\in \Agrp[\vertexV]/\Agrp}(\sum_{\bm{\bm{\chi}}}M_{\bm{\bm{\chi}}})\T^{\bm{a}}\ket{\psi}\otimes \Zp^{\bm{p}_{\bm{\bm{\chi}}}}\ket{\delta_1\bm a}\label{eq:isomorphism-4}\\
    &= \frac{1}{|\Agrp|^{(|\vertexV|-1)/2}}\sum_{\bm a\in \Agrp[\vertexV]/\Agrp}\sum_{\bm{\chi}} \bm{\chi}(\bm a)M_{\bm{\chi}} \T^{\bm a}\ket{\psi}\ket{\delta_1\bm a}\label{eq:isomorphism-5}\\
    &=\frac{1}{|\Agrp|^{(|\vertexV|-1)/2}}\sum_{\bm a\in \Agrp[\vertexV]/\Agrp}\sum_{\bm{\chi}} \bm{\chi}(\bm a)\bm{\chi}(\bm a)^{-1} \T^{\bm a}M_{\bm{\chi}}\ket{\psi}\ket{\delta_1\bm a}\label{eq:isomorphism-6}\\
    &=\frac{1}{|\Agrp|^{(|\vertexV|-1)/2}}\sum_{\bm a\in \Agrp[\vertexV]/\Agrp} \T^{\bm a}(\sum_{\bm{\chi}}M_{\bm{\chi}})\ket{\psi}\ket{\delta_1\bm a}\notag\\
    &=\frac{1}{|\Agrp|^{(|\vertexV|-1)/2}}\sum_{\bm a\in \Agrp[\vertexV]/\Agrp} \T^{\bm a}M \ket{\psi}\ket{\delta_1\bm a} = \Gamma(M\ket{\psi}).\notag
\end{align}
From Eq.~\eqref{eq:isomorphism-4} to Eq.~\eqref{eq:isomorphism-5} we used  $\partial_1 \bm{p}_{\bm{\chi}} = \bm{\chi}$, which implies that
$\Zp^{\bm{p}_{\bm{\bm{\chi}}}}\ket{\delta_1\bm a} = \bm{\chi} (\bm a) \ket{\delta_1 \bm a}$, and from Eq.~\eqref{eq:isomorphism-5} to Eq.~\eqref{eq:isomorphism-6} we used $M_{\bm{\chi}} \T^{\bm{a}} = \bm{\chi}(\bm a)^{-1} \T^{\bm{a}}M_{\bm{\chi}}$.

\end{proof}


From Lemma~\ref{lem:gauging-map-isomorphism} we directly have the following isomorphism between the symmetric code (Def.~\ref{def:symmetric-code}) and the gauged code from Section~\ref{sec:gauging-procedure}.
\begin{corollary}[Isomorphism of gauged code and symmetric code]
\label{cor:sym-gau-isomorphism}
    Let $\T$ be the transversal Clifford symmetry in Def.~\ref{def:transversal-gate}, and assume that the gauging graph $\graphG$ satisfies the graph desiderata in Def.~\ref{def:graph-desiderata}. 
    Let the corresponding gauged code $\cC'$ have stabilizers defined in Thm.~\ref{thm:generic-gauged-stabilizers}.
    Then, the symmetric codespace $\cC_{\rm sym.}$ and the gauged codespace $\cC'$ are isomorphic. The gauging map $\Gamma$ implements this isomorphism.
\end{corollary}
\begin{proof}
Lemma \ref{lem:gauging-map-isomorphism} tells us that $\mathcal H_{\rm sym.}$ is isomorphic to the $+1$ eigenspace of the $\av$ and $\bp$ stabilizers. The extra constraint to arrive at the symmetric codespace and the gauged codespace is to project onto the $+1$ eigenspace of the symmetrized stabilizers and gauged symmetrized stabilizers, respectively. But these are naturally isomorphic through the action of the gauging map through the second statement of Lemma \ref{lem:gauging-map-isomorphism}. 
\end{proof}

We can summarize Lemma \ref{lem:gauging-map-isomorphism} and Corollary \ref{cor:sym-gau-isomorphism} with the following commutative diagram:
\begin{equation}\label{eq:gauging_cube}
\begin{tikzcd}[row sep=3em, column sep=3em, ampersand replacement=\&]
\mathcal H_{\rm sym.}
    \arrow[r, "\Pi_{\cC_{\rm sym.}}"]
    \arrow[d, "\Gamma"', "\cong"]
\& \cC_{\rm sym.}
    \arrow[d, "\Gamma"', "\cong"]
\\
\mathcal H_{\rm gau.}
    \arrow[r, "\Pi_{\cC'}"']
\& \cC'
\end{tikzcd}
\end{equation}

\begin{remark}

We note that the above isomorphisms also apply to the gauged code defined by the unitary $\mathbb{Z}_2$ stabilizers in Theorem \ref{thm:z2-gauged-stabilizers}, since by Theorem \ref{thm:z2-generic-equivalence} this code is equivalent to that obtained in the general case. 
\end{remark}

\subsection{Distance of the gauged code}

Armed with Lemma \ref{lem:gauging-map-isomorphism} we now establish distance preservation in the gauged code.

We define the support of an operator $O$ as $\supp{O}$, the set of physical qudits on which
$O$ is not the identity, and
\[
    \Wt(O):=|\supp{O}|
\]
is its physical weight.

A key property appearing in recent distance proofs for quantum code surgery is the expansion of $\graphG$.
Since we are working with a chain complex on the graph $\graphG$ with coefficients in the finite group $\Agrp$ rather than in $\mathbb{F}_2$, we require a generalized notion of expansion. Intuitively, the Cheeger constant bounds the relationship between the size of a set of vertices (a $0$-chain) and the size of its coboundary (a $1$-chain). 

For an equivalence class $[\bm{g}] = \{\bm{g} + a \bm{1}\}$, let $\tilde{\bm{g}} \in [\bm{g}]$ be the minimal weight representative.


\begin{definition}[Generalized Cheeger constant]
\label{def:generalized-cheeger}
Consider a graph $\graphG$ and coefficient group $\Agrp$. Then the \emph{generalized} Cheeger constant is defined:
\begin{equation}
    \begin{aligned}
    h_{\Agrp}(\graphG):=
    &\min_{\bm g\notin\ker\delta_1}
    \frac{|\supp{\delta_1\bm g}|}{|\supp{\tilde{\bm g}}|}
\end{aligned}.
\end{equation}
\end{definition}




\begin{theorem}[Gauged Code Distance]
\label{thm:space-distance}
Let $d$ be the distance of $\cC$, and let $\nu:=\max_{v\in\vertexV}|\vertexset{v}|$. Let $\Pi_{\cC'}$ denote a projector onto $\cC'$, the gauged code whose stabilizers are defined in Thm.~\ref{thm:generic-gauged-stabilizers} and Def. \ref{def:gauged-code-stabs}. Let $\mathrm{Wt}()$ count the number of matter and gauge qudits on which an operator acts nontrivially. 
Define 
\[
d_g := \min\{\mathrm{Wt}(\mathbb{L}):\Pi_{\cC'} \mathbb{L}\Pi_{\cC'}\not\propto \Pi_{\cC'}\}.
\] 
If the gauging graph satisfies Def.~\ref{def:graph-desiderata}, with generalized Cheeger constant 
\[
h_{\Agrp}(\graphG) \geq 1\]
then
\begin{align}
    d_g\ge d/\nu.
\end{align}
\end{theorem}
\begin{proof}

At a high level, our proof strategy is to break a generic logical $\mathbb{L}$ on the gauged code into its eigen-decomposition on the matter degrees of freedom, and its Pauli support on the gauge degrees of freedom. The weight of the total operator is bounded by the weight of any nontrivial constituent operator, and so we focus on a single component. At this point, we can apply similar arguments to previous work \cite{williamson2024lowoverheadfaulttolerantquantumcomputation, williamson2026fast} to show this component is bounded in weight, namely that it can always be cleaned into the form of a logical on the ungauged code, up to some decorations on the gauging graph. 

Pick an arbitrary $\mathbb{L}$ such that $\Pi_{\cC'}\mathbb{L}\Pi_{\cC'} \not\propto \Pi_{\cC'}$. 
Expand the support of $\mathbb{L}$ on the gauge qudits in term of generalized Paulis, then decompose each matter operator into its eigenoperator decomposition,
    \[
    \mathbb{L} = \sum_{\bm a,\bm b} M_{\bm a,\bm b}\otimes \Xp^{\bm a}\Zp^{\bm b} = \sum_{\bm{\chi}}\sum_{\bm a ,\bm b} (M_{\bm a, \bm b})_{\bm{\chi}} \otimes \Xp^{\bm a} \Zp^{\bm b}= \sum_{i}L_i,\qquad L_i = M_i \otimes \Xp^{\bm a_i}\Zp^{\bm b_i}
    \]
where $\bm a\in \Agrp[\edgeE]$, $\bm b\in \widehat{\Agrp}[\edgeE]$ and $\bm{\chi} \in \widehat{\Agrp}[\vertexV]$, and we relabel $i=(\bm a, \bm b, \bm{\chi})$.
    $M_i$ is a matter eigenoperator with eigenvalue $\bm{\chi}$,    $ \bm a_i\in \Agrp[\edgeE]$ and $ \bm b_i\in \widehat{\Agrp}[\edgeE]$.

    We use the fact that $\mathbb{L}$ is expanded in terms of a tensor product operator basis. On the matter qudits, we have expanded in terms of the eigenoperators of the symmetry action on each vertex set, and thus these basis operators act on at most $\nu$ matter qudits.
    On the edge ancilla, we have expanded in the onsite Pauli basis. 
    The basis for operators over the full system is the tensor product of these local basis operators.
    The support of any non-zero $L_i$ is the identity away from the edge ancilla and matter vertex sets that $\mathbb{L}$ is supported on. Thus, \begin{align}
        \mathrm{Wt}(\mathbb{L})\ge \frac{1}{\nu} \mathrm{Wt}(M_i) + \mathrm{Wt}(\Xp^{\bm{a}_i} \Zp^{\bm{b}_i}) \label{eq:eigendecomp-enlarged-support}
    \end{align}

    for any $L_i$. 
    The factor of $1/\nu$ comes from the fact the eigenoperator expansion may increase the support of $L_i$ in relation to $\mathbb{L}$, but only by the maximum number of sites in a vertex set.

    By linearity of the projector $\Pi_{\cC'}$, if $\Pi_{\cC'}\mathbb{L}\Pi_{\cC'}\not \propto \Pi_{\cC'}$,
    then at least one $L_i$ must satisfy $\Pi_{\cC'}L_i\Pi_{\cC'}\not \propto \Pi_{\cC'}$.

    Separately, we know that every $L_i$ is a simultaneous eigenoperator under conjugation of $\av, \bp$ for $v,p$:
    \[
    \av L_i \av^\dagger = \lambda_{i,v}L_i\qquad \bp L_i \bp^\dagger = \sigma_{i,p}L_i,\qquad\qquad \lambda_{i,v}, \sigma_{i,p}\in U(1).
    \]
    This follows from our choice of operator basis. On matter qudits, $L_i$ consists of a basis operator $M_i$ that satisfies $\T^{\bm c} M_i (\T^{\bm c})^\dagger = \bm{\chi}_i(\bm c) M_i$ for any $\bm c\in \Agrp[\vertexV]$. 
    On gauge qudits, $L_i$ is a generalized Pauli operator. Generalized Pauli operators are eigenoperators under Pauli conjugation, so conjugation by $\bp$ multiplies $L_i$ by a phase. 
    Conjugation by
    $\av$ acts as $\T_v$ on the matter qudits and as a generalized Pauli $\Xp$ on the gauge qubits, each of which multiply $L_i$ by a phase.
        
    Now that we have established that $L_i\ \forall i$ has well-defined eigenvalues under conjugation by $\av,\bp$, we show that any $L_i$ with $\Pi_{\cC'}L_i\Pi_{\cC'}\neq 0$ must satisfy $\lambda_{i,v}=\sigma_{i,p}=1$ for all $v,p$. Since every state in $\cC'$ is stabilized by $\av$ and $\bp$, 
    \[
    \Pi_{\cC'}L_i\Pi_{\cC'} = \Pi_{\cC'}\av L_i\av \Pi_{\cC'}  = \lambda_{i,v}\Pi_{\cC'}L_i\Pi_{\cC'} .
    \]
    If $\lambda_{i,v}\neq 1$, then $\Pi_{\cC'}L_i\Pi_{\cC'} =0$. A similar argument holds for $\bp$: $\sigma_{i,p}\neq 1$ implies that $\Pi_{\cC'}L_i\Pi_{\cC'} =0$.
    
    In summary, $\mathbb{L}$ has at least one $L_i$ such that $\Pi_{\cC'}L_i\Pi_{\cC'}\not \propto \Pi_{\cC'}$, and that $L_i$ must have eigenvalues $\lambda_{i,v}=\sigma_{i,p}=1$ under conjugation by $\av,\bp$ for all $v,p$.
    If this $L_i=M_i\otimes\Xp^{\bm a_i}\Zp^{\bm b_i}$ commutes with all $\bp$, then that means that $\bm a_i\in \ker\delta_2$. Using $H^1(\graphG^\bullet)=0$, this implies that $\ker\delta_2 = \mathop{Im}\delta_1$. Thus, there must exist some $\bm d_i\in \Agrp[\vertexV]$ such that $\bm a_i=\delta_1 \bm d_i$.
    Recall the stabilizer $A(\bm a) = \T^{\bm{a}}\otimes \Xp^{\delta_1\bm a}$ and let its inverse be denoted by $A(\bm a)^{-1} := A(-\bm a) = \T^{-\bm a}\otimes \Xp^{-\delta_1\bm a}$.
    If we multiply $L_i$ by the stabilizer $A(\bm d_i)^{-1}$, then the $\Xp$ support on all gauge qudits cancels out:
    \[
    L_i':=L_iA(\bm d_i)^{-1} = M_i\T^{-\bm d_i}\otimes \Xp^{\bm a_i-\delta_1 \bm d_i}\Zp^{\bm b_i} = M_i\T^{-\bm d_i}\otimes \Zp^{\bm b_i}.
    \]
    Since we only multiplied $L_i$ by a stabilizer $A(\bm d_i)^{-1}$, the logical action of $L_i$ on the codespace is the same as the logical action of $L_i'$ ($\Pi_{\cC'}L_i \Pi_{\cC'} =  \Pi_{\cC'}L_i' \Pi_{\cC'}$), so $L_i'$ has nontrivial action on the gauged code.
    We refer to an $L_i'$ as surviving when $\Pi_{\cC'}L_i' \Pi_{\cC'}\neq 0$.
    
    Next, we show that every surviving $L_i'$ is equivalent, up to stabilizers of $\cC'$, to $\Gamma[N_i]$ for some matter operator $N_i$; in other words, $L_i'$ and $\Gamma[N_i]$ have the same action on $\cC'$.
    Write $L_i' = N_i \otimes \Zp^{\bm b_i}$ where $N_i:= M_i\T^{-\bm d_i}$, and let $\bm{\chi}_i\in \widehat{\Agrp}[\vertexV]$ describe how $N_i$ transforms under the local symmetry, 
    \begin{align}
        \T^{\bm c}N_i (\T^{\bm c})^\dagger     = \bm{\chi}_i(\bm c) N_i,\qquad \forall \bm c\in \Agrp[\vertexV]\notag.
    \end{align}
    By Def.~\ref{def:gauged-operator}, $\Zp(\bm b_i)$ is a valid decoration of $N_i$ if  
    \begin{align}
        \partial_1 \bm b_i = \bm{\chi}_i.
        \label{eq:decorated-requirement}
    \end{align}
    This condition follows from commutation with the $\av$ stabilizers. For $g\in \Agrp$, $\av^g = \T^g_v\otimes \Xp^{\delta_1 \bm e_g}$ where $\bm e_{g,v}$ is $g$ at $v$, and zero elsewhere.
    Conjugating $L_i'$ by this local symmetry gives
    \begin{align}
        \av^gL_i'({\av^g})^\dagger = (\bm{\chi}_i)_v(g)\ (\partial_1\bm b_i)_v(g)^{-1}L_i'.
    \end{align}
    Since a surviving $L_i'$ must commute with every $\av^g$, the phases must cancel for every $v$ and $g$ and therefore $\partial_1\bm b_i = \bm{\chi}_i$.
    Thus, $\Zp^{\bm b_i}$ is a valid decoration of $N_i$.

    The particular decoration $\Zp^{\bm b_i}$ need not be equal to the one chosen in Def.~\ref{def:gauged-operator}.
    But since the two decorations have the same boundary, their difference is a closed $1$-chain.
    Since $H_1(\graphG_\bullet) = 0$, this difference corresponds to a product of $\bp$ stabilizers.
    Therefore, $L_i'$ and $\Gamma[N_i]$ have the same action on $\cC'$, or $\Pi_{\cC'}L_i'\Pi_{\cC'} = \Pi_{\cC'}\Gamma[N_i]\Pi_{\cC'}$.

    Since $L_i'$ has nontrivial action on $\cC$, so does $\Gamma[N_i]$. By the gauging isomorphism (Lemma~\ref{lem:gauging-map-isomorphism}), the action of $\Gamma[N_i]$ on $\cC'$ is the same as the action of $N_i$ on the symmetric code. Since every global symmetry $\T^a$ acts as identity on the symmetric code, $N_i\T^a$ has the same nontrivial action on the symmetric code for every $a\in \Agrp$. 
    Since the symmetric code is a subspace of $\cC$, $N_i\T^a$ cannot act trivially on the original codespace $\cC$. By the definition of the distance $d$ of $\cC$, 
    \begin{align}
        \Wt (N_i\T^a)\ge d\qquad \forall a\in \Agrp.
    \end{align}
    Recall that $L_i=L_i'A(\bm d_i) = N_i\otimes \Zp^{\bm b_i}A(\bm d_i)$.
    Choose $a_i\in \Agrp$ minimizing $|\supp{\bm d_i - a_i \id}|$.
    This is the reduced support of $\bm d_i$ that appears in the definition of the generalized Cheeger constant.
    The matter support $M_i = N_i\T^{\bm d_i}$ can differ from that of $N_i\T^{a_i}$ by at most $\nu |\supp{\bm d_i - a_i \id}|$, while the gauge support has at least $|\supp{\delta_1\bm d_i}|$. 
    Thus,
    \begin{align}
        \Wt(M_i) \ge \Wt(N_i \T^{a_i}) - \nu |\supp{\bm d_i - a_i \id}|
    \end{align}
    and on the gauge qudits, 
    \begin{align}
        \Wt\!\left(\Xp^{\bm a_i}\Zp^{\bm b_i}\right) \ge |\supp{\delta_1\bm d_i}| \ge |\supp{\bm d_i-a_i\id}|,
    \end{align}
    The second inequality uses the expansion condition that $h_{\Agrp}(\graphG) \ge 1$.

    Finally, the eigenoperator decomposition can enlarge matter support by at most a factor of $\nu$ (see Eq.~\eqref{eq:eigendecomp-enlarged-support}).
    Thus,
    \begin{align}
        \Wt(\mathbb{L}) &\ge \frac{1}{\nu}
        \Wt(M_i)
        + \Wt \!\left(\Xp^{\bm a_i}\Zp^{\bm b_i}\right)\\
        &\ge \frac{1}{\nu}
        \left(\Wt(N_i \T^{a_i})
        - \nu |\supp{\bm d_i -a_i \id }|\right) + 
        |\supp{\bm d_i -a_i \id}|\\
        &\ge \frac{d}{\nu}.
    \end{align}

    
\end{proof}

\begin{remark}[The $1/\nu$ factor in the distance lower-bound]
    The factor of $1/\nu$ in the bound above $d_\mathrm{g}\ge d/\nu$ arises from the eigenoperator decomposition and is specific to our argument for gauging a Clifford logical operator.
    Let $L$ be an operator that acts nontrivially on the gauged code. 
    After taking the eigenoperator decomposition of $L$, the rest of the proof shows that at least one term must have nontrivial action and has weight at least the original code distance $d$.
    However, for a general Clifford symmetry, using the eigenoperator decomposition on $L$ can enlarge the matter support within each vertex set.
    Since a vertex set contains at most $\nu$ matter qudits, the weight of the selected term can be as large as $\nu\Wt(L)$. Consequently, $\nu\Wt(L)\ge d$, leading to the distance bound shown.

    \vspace{6pt}
    The $1/\nu$ factor does not arise for Pauli surgery.
    This is because Pauli operators already have definite eigenvalues under conjugation by a Pauli symmetry, and hence the operator expansion does not enlarge physical qudit support. In this case, if we require $h_{\Agrp}(\graphG) \geq \nu $, the argument above leads to a gauged code distance satisfying $d_\mathrm{g}\ge d$, even if a vertex set contains multiple matter qudits. This is similarly true for surgery on any operator that is written in a definite charge basis on each vertex.
\end{remark}

\begin{remark}
    For some symmetries such as the transversal CNOT, a minimal vertex set contains one qubit from each codeblock, and thus $\nu=2$. 
    Nevertheless, we expect the distance of the gauged code to be preserved ($d_\mathrm{g}\ge d$). This is because even though the eigenoperator expansion extends support between codeblocks, the support on any single codeblock does not grow. We do not establish a proof of this claim here, we expect such a proof would require careful tracking of support between the two separate codeblocks.
\end{remark}

The following corollary is useful in the proof of the fault distance for the full spacetime gauging protocol in the next section. It demonstrates that we are free to add relative phases when we are breaking up the symmetry $\T$. When we perform the gauging circuit in practice, we start in a random configuration of $\av$ measurement outcomes. These outcomes may not correspond to $\T = +1$; the following corollary ensures that the resulting projected state still has distance $d_g\geq d/\nu$.

\begin{corollary}[Distance in the presence of random $\av$] 
    The codespace defined by $\T = \omega^n$ for $n = 0, \dots, p-1$ and a random set of $\av$ eigenvalues $a_v$ such that $\prod_v a_v = \omega^n$ has distance $d_g \geq d/\nu$, if the graph desiderata (Def. \ref{def:graph-desiderata}) are met.
\end{corollary}
\begin{proof}
    This follows from a re-definition of the symmetry $\T$ and the $\av$ stabilizers: 
    \[
    \T \to \omega^{-n} \T, \quad \av \to a_v^* \av
    \]
    Now we have $\T = \prod_v \av = +1$. Define the gauged code in terms of the redefined symmetry operator and $\av$ terms. Apply Theorem \ref{thm:space-distance}.
\end{proof}

\ifdefined\ismainfile
\else

\bibliography{bib}
    \end{document}
\fi
